\begin{figure*}

\begin{minipage}{0.5\linewidth}
\centering
\begin{subfigure}{0.49\linewidth}
\centering
\includegraphics[width=\linewidth]{../binder/binder-2026-08-29-pcomplexity-gee-0s.ipynb/binder/teeplots/2026-08-29-pcomplexity-gee-0s/other=co-evo+viz=subplots+what=self-adaptation+ext=.pdf}

\vspace{-2ex}

\caption{\footnotesize baseline co-evo bkgd. \@ x0's}
\label{fig:reg-self-adaptation:co-evo-0s}
\end{subfigure}
\begin{subfigure}{0.49\linewidth}
\centering
\includegraphics[width=\linewidth]{../binder/binder-2026-08-29-pcomplexity-gee-0s.ipynb/binder/teeplots/2026-08-29-pcomplexity-gee-0s/other=none+viz=subplots+what=self-adaptation+ext=.pdf}

\vspace{-2ex}

\caption{\footnotesize baseline no bkgd. \@ x0's}
\label{fig:reg-self-adaptation:none-0s}
\end{subfigure}

\begin{subfigure}{0.49\linewidth}
\centering
\includegraphics[width=\linewidth]{../binder/binder-2026-07-06-pcomplexity-gee.ipynb/binder/teeplots/2026-07-06-pcomplexity-gee/other=co-evo+viz=subplots+what=self-adaptation+ext=.pdf}

\vspace{-2ex}

\caption{\footnotesize baseline co-evo bkgd. \@ x9's}
\label{fig:reg-self-adaptation:co-evo-9s}
\end{subfigure}
\begin{subfigure}{0.49\linewidth}
\centering
\includegraphics[width=\linewidth]{../binder/binder-2026-07-06-pcomplexity-gee.ipynb/binder/teeplots/2026-07-06-pcomplexity-gee/other=none+viz=subplots+what=self-adaptation+ext=.pdf}

\vspace{-2ex}

\caption{\footnotesize baseline no bkgd. \@ x9's}
\label{fig:reg-self-adaptation:none-9s}
\end{subfigure}
\end{minipage}%
\begin{minipage}{0.5\linewidth}
\caption{%
\textbf{Correlation between self-adaptation and phenotype complexity across 100-stint replicates.}
\footnotesize
Pearson and Spearman correlation statistics annotated.
Shaded bands indicate bootstrap 95\% confidence for OLS regression.
Fitness values represent proportion ``measuring stick'' competitions won against comparator strains.
Self-adaptation, then, is the difference between fitness in self-dense biotic context and fitness in baseline context.
Top panels show results averaged per replicate for exploratory assays across stints 10, 20, 30, 40, 50, 60, 70, 80, 90, 100; bottom panels show results for confirmatory assays at stints 9, 19, 29, 39, 49, 59, 69, 79, 89, 99 (see note in Supplementary Figure \ref{fig:complexity-fitness-correlation-genetic-9s}).
Similar patterns arise whether baseline context is either co-evolved background (panels \labelcref{fig:reg-self-adaptation:co-evo-0s,fig:reg-self-adaptation:co-evo-9s}) or no background (panels \labelcref{fig:reg-self-adaptation:none-0s,fig:reg-self-adaptation:none-9s});
note, though, that Spearman correlation under no-background baseline at x9's does not reach significance at $\alpha = 0.05$.
Samples drawn from original 40-replicate/100-stint trials, which include the case study replicate.
}
\label{fig:reg-self-adaptation}
\end{minipage}

\end{figure*}
